\documentclass[12pt]{article}

\usepackage[margin=0.75in]{geometry}
\usepackage{setspace}
\usepackage{parskip}
\usepackage{hyperref}
\usepackage{titlesec}
\usepackage{xcolor}
\usepackage{array}
\usepackage{ragged2e}

\definecolor{PrimaryRed}{HTML}{8B1E1E}
\definecolor{SoftBlue}{HTML}{DDEEFF}
\definecolor{DarkGray}{HTML}{333333}

\hypersetup{
    colorlinks=true,
    linkcolor=black,
    urlcolor=black,
    citecolor=black
}

\titleformat{\section}{\large\bfseries\color{PrimaryRed}}{}{0em}{}
\titleformat{\subsection}{\normalsize\bfseries}{\thesubsection}{0.5em}{}

\setstretch{1.15}

\begin{document}

\begin{titlepage}
\newgeometry{margin=0.85in}
\thispagestyle{empty}

\begin{center}
    \vspace*{0.25in}
    {\large University of Rochester\\}
    {\normalsize Warner School of Education\\[0.35in]}

    {\color{PrimaryRed}\rule{\textwidth}{1.4pt}}\\[0.22in]
    {\Huge\bfseries Joy-Oriented Literacy Pursuit\\[0.12in]}
    {\Large Formal Reflective Journal --- Module 3 Update\\[0.22in]}
    {\color{PrimaryRed}\rule{0.72\textwidth}{0.7pt}}\\[0.32in]

    \colorbox{SoftBlue}{%
        \parbox{0.86\textwidth}{%
            \centering\vspace{0.12in}
            \textbf{Digital and Media Literacy as Joy, Access, and Identity}\\[0.05in]
            Updated with Module 3: Week 5 readings on digital literacy, criticality, and culturally responsive design
            \vspace{0.12in}
        }%
    }
\end{center}

\vspace{0.75in}

\begin{center}
\renewcommand{\arraystretch}{1.35}
\begin{tabular}{>{\bfseries\RaggedLeft}p{1.55in} p{4.85in}}
Student & Piter Garcia \\
Course & EDU498: Culturally-Historically Responsive Literacies \\
Class Session & Class 5 (Module 3 Update) \\
Assignment & Joy-Oriented Literacy Pursuit \\
Context & Teaching Placement Reflection \\
Submission Date & June 2026 \\
\end{tabular}
\end{center}

\vfill

\begin{center}
\begin{minipage}{0.78\textwidth}
\centering
\itshape
Digital literacy is not a technology topic added to real learning. It is real learning --- the way students already speak, build, represent, and move through the world.
\end{minipage}
\end{center}

\restoregeometry
\end{titlepage}
\clearpage

\section*{What I Created and Who I Shared It With}

For my joy-oriented literacy pursuit, I created and facilitated classroom-based technology and design activities where students could experience literacy through building, collaboration, speaking, movement, problem-solving, and shared excitement. Instead of treating literacy only as reading or writing, I focused on moments where students used language, materials, technology, and relationships to make meaning. The main examples from my teaching placement were the Kindergarten circuit and fan lesson, the Grade 4 Minecraft activity, and the Grade 5 airplane and fan-circuit challenge. These activities were shared with students in my placement classrooms, but they were also shared with the classroom community itself, because students built, watched, helped, responded, and learned from one another in real time.

Module 3 deepened my understanding of why those moments matter. Knobel and Lankshear name practices like the ones I witnessed as \textit{new literacies}: meaningful, social, multimodal ways of making meaning that schools often fail to recognize as literacy. Price-Dennis, Holmes, and Smith extend this by arguing that inclusive digital literacy instruction requires teachers to see students' existing digital practices as resources, not distractions. What I witnessed in Grade 4 Minecraft and Grade 5 circuit design were not students using technology; they were students \textit{reading, writing, and communicating} through digital and physical tools --- in exactly the sense the Module 3 readings describe.

\section*{My Creative and Learning Process}

My creative process began with a question: how can technology lessons become more than task completion? I wanted students to feel that literacy was something they could touch, build, speak, test, revise, and share. Instead of designing lessons where I explained everything and students simply copied the steps, I looked for ways to make students active contributors to the learning process.

Module 3 added a critical dimension to that process. Brooks and Frankel's work on transnational digital literacy helped me understand that students may already be using technology to move across languages, countries, and communities --- long before they arrive at any lesson I design. That means the classroom is not introducing students to digital literacy; it is meeting students where their digital literacies already exist. For students who are Dominican, multilingual, neurodivergent, or navigating multiple identities, digital spaces can already be places of voice, community, and belonging. Stornaiuolo and Thomas reinforce this: young people use digital tools to speak, organize, and participate in public life. When a Grade 5 student explains their airplane circuit to a peer, or when a Grade 4 student stays to help a classmate complete a Minecraft build, those are not small moments. They are exactly the kind of participatory digital literacy Stornaiuolo and Thomas describe.

Pinkard's ecosystem model also changed how I understand my role as a designer. Digital literacy does not happen when a device is added to a task. It happens when students have tools, mentors, audiences, feedback, and real reasons to create. In my placement, that ecosystem was present when students had clear roles, visible materials, a peer support structure, and the experience of shared success. That was not accidental. It was designed --- and Module 3 gave me language to name and defend that design.

\section*{How This Cultivated Joy, Love, and Aesthetic Fulfillment}

This literacy pursuit cultivated joy because students were able to see their learning become real. In the Kindergarten circuit lesson, the joy came from collective anticipation: students named parts, held materials, counted down together, and watched the fan turn on. That moment became a shared literacy event. The joy was not only in the fan working; it was in the class experiencing the success together.

Module 3 added another layer to how I understand that joy. Serpagli and Mensah show that social media-style practices --- documenting, sharing, explaining, reacting --- can support learning when students use familiar digital forms to participate. The countdown in the Kindergarten lesson functioned like that: it gave students a participatory format that made the science visible, social, and memorable. The Grade 5 airplane teams did the same thing through a different channel --- building, testing, adjusting, naming, and demonstrating. These were not just activities. They were, in the terms of Knobel and Lankshear, \textit{new literacies in action}.

The pursuit also cultivated love through relationships. In Grade 4 Minecraft, students who already understood the task helped classmates who were still working. Module 3 helped me name that moment through Pinkard's ecosystem lens: peer mentors, audience, and community feedback are not optional additions to digital literacy. They are what makes digital literacy real and sustaining. Aesthetic fulfillment in this context means students experiencing their own competence as something beautiful --- something worth sharing, worth teaching, worth celebrating.

\section*{Muhammad's Literacy Pursuit and Module 3}

This project connects to Muhammad's idea of literacy pursuit because the activities were not designed only to teach isolated skills. They created opportunities for students to build identity, skills, intellect, criticality, and joy. Module 3 deepens every one of those dimensions through digital and media literacy.

Criticality, in particular, gained new weight from the Module 3 readings. Melo-Pfeifer and Gertz argue that students need critical intercultural news literacy: the capacity to question headlines, images, missing voices, and repeated assumptions. The Common Sense Digital Citizenship materials add confirmation bias, online identity, privacy, audience awareness, and upstander communication. These are not separate from joy-oriented literacy. They are part of how students learn to move through digital spaces with agency and dignity --- to be creators and critics, not just consumers. Muhammad's framework supports this: criticality and joy are not opposites. A student who knows how to question a media message, advocate for their community online, or recognize when their identity is misrepresented is a student who has been given more tools to pursue joy on their own terms.

For students who are Black, Latino, immigrant, multilingual, disabled, or otherwise outside the assumed ``default learner,'' digital spaces carry real risk alongside real possibility. Teaching digital literacy without teaching criticality leaves those students exposed. Teaching criticality without joy leaves them exhausted. Module 3 argues for both --- and so does Muhammad.

\section*{How This Informs My Future Teaching}

This experience, now deepened by Module 3, will inform my teaching by reminding me that digital literacy is identity work. Students do not enter digital spaces as blank users. They bring language, culture, family, community, interests, histories, bodies, and lived experiences. Brooks and Frankel name this explicitly for transnational students; I carry it forward for all students whose identities are already present in how they use, question, and create in digital spaces.

In practice, I would continue designing projects where students create something meaningful and share it with an authentic audience. Students could build Minecraft models connected to a science concept, create bilingual or multimodal explanations of a design challenge, make collaborative classroom murals, create technology tutorials for peers, document a circuit build through a video explanation, or produce short critical media responses to news articles affecting their communities. I would also treat peer teaching as a literacy practice in itself --- one that builds community responsibility alongside technical skill.

My biggest takeaway from Module 3 is that digital tools can open doors \textit{or} reinforce walls, depending on how teachers design instruction. If I add a device to a traditional task without changing the structure, I have not created access --- I have just digitized the same narrow expectations. Strong digital literacy instruction gives students multiple ways to respond, teaches them to evaluate media carefully, supports responsible participation, and treats their existing digital practices as starting points for learning, not obstacles to overcome. That is what joy-oriented, culturally and historically responsive digital literacy looks like in practice.

\section*{Evidence Notes}

This updated reflection draws on Module 3 sources (Knobel \& Lankshear; Price-Dennis, Holmes, \& Smith; Stornaiuolo \& Thomas; Brooks \& Frankel; Serpagli \& Mensah; Melo-Pfeifer \& Gertz; Common Sense Digital Citizenship; Pinkard) alongside the original placement evidence from the Kindergarten circuit and fan lesson, the Grade 4 Minecraft peer-support activity, and the Grade 5 airplane and fan-circuit challenge.

\end{document}
